\begin{frame}{Motion, loss, and quasi-stationarity}
\small
The underlying motion is stochastic, but retention introduces loss:
\[
 P\mathbf{1}=\mathbf{1},\qquad
 A=P\operatorname{diag}(r),\qquad
 0\le r_i\le1,\qquad A\mathbf{1}\le\mathbf{1}.
\]
The stationary distribution of the motion satisfies
\[
 \pi^TP=\pi^T,\qquad \pi^T\mathbf{1}=1.
\]
For irreducible $A$, its normalized left Perron vector satisfies
\[
 q^TA=\Lambda q^T,\qquad q^T\mathbf{1}=1,
 \qquad \boxed{q^TA\mathbf{1}=\Lambda}.
\]
If $\Lambda>0$, conditioning on survival gives
\[
 \frac{q^TA}{q^TA\mathbf{1}}=q^T.
\]
Thus $q$ is a \emph{quasi-stationary distribution}, and $\Lambda$
is its one-step survival probability. At uniform retention $r=K\mathbf{1}$,
\[
 A=KP,\qquad q=\pi,\qquad \Lambda=K.
\]
\end{frame}

%================================================


\begin{frame}{Random-matrix perturbation on the actual graph}
\small

For the fixed graph, write
\[
    r=K\mathbf{1}+\varepsilon,
    \qquad
    A=KP+P\operatorname{diag}(\varepsilon).
\]

Here $P\mathbf{1}=\mathbf{1}$: $P$ describes motion before loss.
Let $\pi$ be the stationary distribution of $P$, and define
\[
    \pi^TP=\pi^T,
    \qquad
    \pi^T\mathbf{1}=1,
    \qquad
    \Pi=\mathbf{1}\pi^T,
\]
\[
    H=(I-P+\Pi)^{-1}(P-\Pi).
\]

A second-order expansion of the Perron root gives
\[
    \boxed{
    \Lambda\approx
    K+\pi^T\varepsilon+
    \frac{1}{K}
    \pi^T\operatorname{diag}(\varepsilon)H\varepsilon.
    }
\]

\textbf{Assumptions:}
$K>0$, irreducible $P$, and perturbations small enough
relative to the reduced resolvent scale.

\medskip

{\footnotesize
This is a local expansion, not a universal large-matrix theorem.
Non-normality and very small neighborhoods can affect its accuracy.
We therefore compare it with actual simulated eigenvalues.
}
\end{frame}


%================================================

\begin{frame}{Deriving the second-order correction}
\small

Put $D=\operatorname{diag}(\varepsilon)$ and
$A(t)=KP+tPD$. Expand
\[
    v(t)=\mathbf{1}+tu+O(t^2),
    \qquad
    \pi^Tu=0,
\]
\[
    \lambda(t)=K+ta+t^2b+O(t^3).
\]

The first-order eigenvector equation gives
\[
    K(I-P)u=P\varepsilon-a\mathbf{1}.
\]
Multiplying by $\pi^T$ yields
\[
    a=\pi^T\varepsilon,
    \qquad
    u=\frac{H\varepsilon}{K}.
\]

At second order, multiplication by $\pi^T$ gives
\[
    b=\pi^TPDu
     =\pi^TDu
     =\frac{1}{K}
       \pi^T\operatorname{diag}(\varepsilon)H\varepsilon.
\]

Thus
\[
    \Lambda\approx K+a+b.
\]

{\footnotesize
The approximation evaluates a local expansion about $t=0$
at $t=1$. Its finite-sample accuracy must be checked.
}
\end{frame}


%================================================

\begin{frame}{Predicting the mean and distribution}
\small

Under fixed-count resident assignment, the covariance is exactly
\[
    C=\operatorname{Cov}(\varepsilon)
    =
    \frac{K(1-K)}{N-1}
    \left[
        N\operatorname{diag}(1/n_i)
        -\mathbf{1}\mathbf{1}^T
    \right].
\]

Since $\mathbb{E}[\varepsilon]=0$,
\[
    \boxed{
        m_0\approx
        K+\frac{1}{K}
        \sum_{i,j}\pi_iH_{ij}C_{ij}
    },
\]
\[
    \boxed{
        \operatorname{Var}(\Lambda)\approx\pi^TC\pi
    }.
\]

If the linear term dominates and a finite-population
central limit theorem applies,
\[
    \Lambda\ \dot\sim\
    \mathcal{N}\!\left(
        K+\frac{1}{K}\sum_{i,j}\pi_iH_{ij}C_{ij},
        \ \pi^TC\pi
    \right).
\]

{\footnotesize
The mean correction is second order; the variance is first order.
The normal approximation requires that no resident contribution
dominate the weighted fluctuation.
If $\pi_i=n_i/N$, the linear variance vanishes.
For complete mixing, $H=0$ and $\Lambda=K$ exactly.
}
\end{frame}


%================================================
